\begin{abcliste}
\abc Wir behandeln die Platte als schwarzen Strahler ($\varepsilon = 1$); eine reale, dunkle Herdplatte strahlt etwas weniger ab. Mit $T = \TP$ und der Fläche $A = \pi d^2/4$ ist die in der Zeit $t$ abgestrahlte Energie (Stefan-Boltzmann, FoTaBe S.~172)
\begin{align}
 E &= \sigma\,A\,T^4\,t = \EPF \\
   &= \sig\cdot\frac{\pi\cdot\left(\dP\right)^2}{4}\cdot\left(\TP\right)^4\cdot\tS \\
   &= \EP \approx \result{\EPP{3}{2}}
\end{align}
Die Strahlung, welche die Platte aus der Umgebung aufnimmt, ist dabei nicht berücksichtigt.

\abc Nach dem Wien'schen Verschiebungsgesetz (FoTaBe S.~182) ist
\begin{align}
 \lambda_\text{max} &= \lPF = \frac{\bW}{\TP} = \lP \approx \result{\lPP{-6}{2}}
\end{align}
Das liegt im Infrarot; die Platte glüht bei dieser Temperatur noch nicht sichtbar.
\end{abcliste}
